\section{Model, data, and experiment designs}\label{sec:method}

\subsection{Model and data}

All results come from PolicyEngine-US version \genModelVersion{}
\citep{policyengineus}, an open-source static microsimulation of United
States federal and state tax and benefit law, executed through the
\texttt{policyengine.py} runtime (version \genRunnerVersion{}), which
certifies the pairing of model version and dataset build. The 2030 baseline
is current law as of that model version, including the 2025 reconciliation
act. The dataset is \texttt{populace-us} build \texttt{\genDataBuildShort}
(\genDataBuildDate{}) \citep{populaceus}, an openly published enhanced-CPS
microdata file for 2024, uprated by the model to the analysis years; scenario
growth is defined in excess of the CBO baseline path, which is exactly right
only where the model's own uprating tracks CBO's projection. The mechanism
sweep of Section~\ref{sec:sweep} is run at \genSweepYear{}; the calibrated
scenarios are run at 2030, matching the Budget Lab's horizon. In the 2030
baseline, SPM poverty is \genBaselinePovertyPct\% and the household
net-income Gini is \genBaselineGini{} (household-level, unequivalized;
Section~\ref{sec:discussion} discusses comparability).

The model covers federal individual income tax (including capital gains
rates, the net investment income tax and the alternative minimum tax),
payroll taxes on both employee and employer sides plus self-employment tax,
refundable credits (EITC, refundable CTC and others), the major means-tested
transfer programs (SNAP, SSI, TANF, WIC, school meals, housing subsidies,
broadband subsidies), and the tax and benefit codes of all fifty states and
DC. Health programs (Medicaid, CHIP, ACA subsidies) are valued in the model
but excluded from net-income accounting in these runs by the model's default
simulation switch, so no result here contains a health-program response in
either direction. Head Start and Early Head Start are likewise valued in the
model but kept out of net income in the calibrated scenarios, as current
PolicyEngine does. The mechanism sweep predates that correction and still
counts both programs in benefit outlays and net income (see the correction
before the introduction). Everything is static: no labor supply, behavioral,
avoidance or enforcement responses, matching the Budget Lab's individual
microsimulation.

\textbf{Accounting.} Net government revenue is defined as taxes minus
transfers: household tax before refundable credits (federal and state,
including employee-side payroll) plus employer payroll tax, minus refundable
credits and benefit outlays. An identity check --- the change in market
income minus the change in household net income must equal the household-side
revenue change --- closes to within \$0.001B in every scenario (maximum
absolute residual \$\genIdentityMaxResidual B at three decimals).
Where the paper compares against the Budget Lab
(Section~\ref{sec:reconciliation}), it restates results on their narrower
basis: federal individual income tax net of refundable credits, plus payroll.

\subsection{Design I: the mechanism sweep}\label{sec:method-sweep}

The sweep isolates composition. In 10-point steps from 0\% to 100\%, it
reduces every person's positive employment and self-employment income by the
step share and redistributes the same weighted total to positive capital
income in proportion to existing holdings, holding modelled market income
fixed. Capital income here means the six directly held flows (long- and
short-term capital gains, taxable interest, qualified and non-qualified
dividends, and rental income); the calibrated scenarios of
Section~\ref{sec:method-scenarios} use a wider nine-variable set that adds
pass-through and taxable retirement income, which is why the two baselines'
capital aggregates differ. The \genSweepYear{} sweep baseline contains
\$\genSweepLaborT T of positive labor income and \$\genSweepCapitalT T of
positive capital income on the six-variable definition, with
\genSweepHHCapSharePct\% of households receiving any positive capital income
--- so the sweep concentrates income by construction, and the question is
what current law does about it.

The sweep's limitation is deliberate: a pure reallocation holds output fixed,
while forecasters expect AI to raise output \emph{and} tilt its composition.
The sweep therefore answers ``what does composition alone do to current
law?'', and the calibrated scenarios answer ``what happens on the paths
forecasters actually expect?''. The two connect through a
forecast-equivalence mapping: the proportional reduction in labor's income
share that each tilt produces, $k = 1 - (1+g_L)/(1+g_Y)$, is also the sweep
position delivering the same share reduction, placing the scenarios at
\genForecastShiftSlowPct\% (Slow), \genForecastShiftModPct\% (Moderate) and
\genForecastShiftRapidPct\% (Rapid) --- the sweep's forecast-relevant range
is its first decade, and the rest is out-of-sample stress test.

\subsection{Design II: forecaster-calibrated scenarios}\label{sec:method-scenarios}

The scenarios adopt the Budget Lab's calibration of \citet{karger2026}
exactly. Each scenario is defined by an annual real GDP growth rate under AI
(Slow 2.0\%, Moderate 2.6\%, Rapid 3.3\%) and a 2030 labor share of factor
income (55.0\%, 53.8\%, 51.3\%). Growth is expressed in excess of the CBO
baseline path: the five-year output shock is
$g_Y = (1+r_{\mathrm{ai}})^5 / G_{\mathrm{CBO}} - 1$, and the factor shocks
follow from the target shares,
\begin{equation}
g_K = \frac{\theta^1_K (1+g_Y) - \theta^0_K}{\theta^0_K}, \qquad
g_L = \frac{\theta^1_L (1+g_Y) - \theta^0_L}{\theta^0_L}.
\end{equation}
Two inputs their report uses but does not print --- the cumulative CBO
baseline factor $G_{\mathrm{CBO}} = 1.0976$ and the pre-shock labor share
$\theta^0_L = 0.5550$ --- are recovered by inverting their published
formulas against the nine derived values in their Table~1; unit tests
reproduce all nine to within 0.005\,pp and re-derive the labor share by grid
search so the constants cannot drift.

The wage-inequality channel is their stylization $\lambda = 1 \pm g_Y$: a
spread transform on log labor income,
$\ln Y^1_i = \mu^1 + \lambda\,(\ln Y^0_i - \mu^0)$, equivalently
$Y^1_i = C\,Y_i^{0\,\lambda}$ with $C$ in closed form, with the aggregate
pinned to $(1+g_L)$ times baseline. Compressive ($\lambda<1$) narrows the
wage distribution around the same aggregate; expansive ($\lambda>1$) widens
it; proportional scales every wage equally. The transform is applied at the
person level because the payroll result turns on a per-person threshold, the
Social Security taxable maximum (\$\genSSCap{} in 2030; at baseline
\genShareWagesAboveCapPct\% of wage income and
\genShareWorkersAboveCapPct\% of workers sit above it).

The capital shock scales nine capital income variables --- long- and
short-term capital gains, taxable interest, qualified and non-qualified
dividends, rental income, partnership and S-corporation income, taxable
pension income and taxable IRA distributions --- in proportion to existing
holdings, after routing pass-through income through the Budget Lab's
labor/capital split rule. The wage-spread transform is not applied to the
labor portion of pass-through income, which follows the pass-through rule
instead. Appendix~\ref{app:routing} traces every shocked
variable into IRS gross income and household market income, confirming the
shock reaches the tax base rather than being recomputed away; the labor
aggregate lands on its target to within relative error
\genLaborTargetMaxErr{}. Each scenario also has a shares-fixed
counterfactual growing both factors at $g_Y$, the comparison behind the
Budget Lab's ``roughly twice as large'' finding.

Two further axes vary definitional assumptions their report flags. A
\emph{realization rate} routes a fraction of the incremental capital flow
outside the taxable base (their assumption corresponds to 1.0); and a
\emph{capital-scope} variant excludes taxable pension income and IRA
distributions from the capital set, since those are taxed at ordinary rates
and their inclusion works against the ``capital is taxed lightly''
mechanism.
